\documentclass{article}
\usepackage{geometry}
\geometry{margin=1.5cm, vmargin={0pt,1cm}}
\setlength{\topmargin}{-1cm}
\setlength{\headheight}{12.5pt}
\setlength{\paperheight}{29.7cm}
\setlength{\textheight}{25.3cm}

% useful packages.
% \usepackage[UTF8]{ctex}
\usepackage{fancyhdr}
\usepackage{layout}

% import common macros
% % % % % % % % % % % % % % % % % % % % % % % % % % % % % % % %
% Common Macros for LaTeX
% % % % % % % % % % % % % % % % % % % % % % % % % % % % % % % %

% Useful packages
\usepackage{amsfonts}
\usepackage{amsmath}
\usepackage{amssymb}
\usepackage{amsthm}
\usepackage{enumerate}
\usepackage{graphicx}
\usepackage{multicol}
\usepackage[ruled,linesnumbered]{algorithm2e}

% Math operators and common symbols
\newcommand{\dif}{\mathrm{d}}
\newcommand{\innerProd}[1]{\left\langle #1 \right\rangle}
\newcommand{\difFrac}[2]{\frac{\dif #1}{\dif #2}}
\newcommand{\pdfFrac}[2]{\frac{\partial #1}{\partial #2}}
\newcommand{\T}{\mathrm{T}}
\newcommand{\norm}[1]{\left\| #1 \right\|}
\newcommand{\abs}[1]{\left| #1 \right|}

% Vector and Matrix shortcuts
\newcommand{\vecTwo}[2]{
  \begin{bmatrix}
    #1 \\
    #2
\end{bmatrix}}
\newcommand{\vecThree}[3]{
  \begin{bmatrix}
    #1 \\
    #2 \\
    #3
\end{bmatrix}}
\newcommand{\vecTwoH}[2]{
  \begin{bmatrix}
    #1 & #2
  \end{bmatrix}
}
\newcommand{\vecThreeH}[3]{
  \begin{bmatrix}
    #1 & #2 & #3
  \end{bmatrix}
}

% Common fields
\newcommand{\R}{\mathbb{R}}
\newcommand{\C}{\mathbb{C}}
\newcommand{\Z}{\mathbb{Z}}
\newcommand{\N}{\mathbb{N}}

% Solutions and proofs
\newenvironment{solution}{
\begin{proof}[解]}{
\end{proof}}

% Utility to get environment variables (requires lualatex)
\newcommand{\getEnv}[2]{\directlua{tex.print(os.getenv("#1") or "#2")}}


% Name and uID
\newcommand{\name}{\getEnv{NAME}{NAME}}
\newcommand{\uid}{\getEnv{UID}{UID}}
\newcommand{\course}{Point Set Topology}
\newcommand{\hwNum}{9}

\begin{document}

\pagestyle{fancy}
\fancyhead{}
\lhead{\zhstr{\name} (\uid)}
\chead{\course\ \#\hwNum}
\rhead{\today}

\section{Question 1}

Show that if $X_i$ is separable for each $i\in\N$, then the product space
\[
  \left(\prod_{i\in\N}X_i,\mathcal{T}_p\right)
\]
is separable.

\begin{solution}
  If some $X_i$ is empty, then the product is empty and is trivially
  separable. Hence assume every $X_i$ is nonempty. For each $i\in\N$,
  choose a countable dense set $D_i\subseteq X_i$ and fix $a_i\in D_i$.

  Let $D$ be the set of sequences which agree with $(a_i)_{i\in\N}$ at
  all but finitely many coordinates and whose remaining coordinates belong
  to the corresponding $D_i$:
  \[
    D=\left\{x=(x_i)\in\prod_{i\in\N}X_i:
      x_i\in D_i\ \text{for all }i,\quad
    \{i:x_i\ne a_i\}\ \text{is finite}\right\}.
  \]
  This set is countable. Indeed, for each finite $F\subseteq\N$, the set
  of elements of $D$ which can differ from $(a_i)$ only on $F$ is a finite
  product of countable sets. There are only countably many finite subsets
  of $\N$, so $D$ is a countable union of countable sets.

  It remains to show that $D$ is dense. Let
  \[
    W=\prod_{i\in\N}U_i
  \]
  be a nonempty basic open set in the product topology. There is a finite
  set $F\subseteq\N$ such that $U_i=X_i$ for every $i\notin F$. For each
  $i\in F$, density of $D_i$ gives $d_i\in D_i\cap U_i$. Define
  \[
    d_i'=
    \begin{cases}
      d_i,&i\in F,\\
      a_i,&i\notin F.
    \end{cases}
  \]
  Then $d'=(d_i')\in D\cap W$. Thus every nonempty basic open set meets
  $D$, so $D$ is dense and the product is separable.
\end{solution}

\section{Question 2}

Let $(X,d)$ be a metric space. Show that the following properties are
equivalent:
\begin{enumerate}[(i)]
  \item $X$ is second countable;
  \item $X$ is Lindel\"of;
  \item $X$ is separable.
\end{enumerate}

\begin{solution}
  We prove that each condition is equivalent to (i).

  \textbf{(i) $\Rightarrow$ (ii).}
  Let $\mathcal{B}$ be a countable base for $X$, and let $\mathcal{U}$ be
  an open cover of $X$. For every $B\in\mathcal{B}$ which is contained in
  at least one member of $\mathcal{U}$, choose one such member $U_B$.
  The collection of the chosen $U_B$'s is countable. It covers $X$:
  if $x\in X$, choose $U\in\mathcal{U}$ with $x\in U$, then choose
  $B\in\mathcal{B}$ with $x\in B\subseteq U$. Thus $x\in U_B$.
  Therefore $X$ is Lindel\"of.

  \textbf{(ii) $\Rightarrow$ (i).}
  For each $n\in\N$, the family
  \[
    \{B(x,1/n):x\in X\}
  \]
  is an open cover of $X$. By the Lindel\"of property, choose a countable
  subcover
  \[
    \{B(x_{n,k},1/n):k\in\N\}.
  \]
  Set
  \[
    \mathcal{B}=\{B(x_{n,k},1/n):n,k\in\N\}.
  \]
  This is countable. To see that it is a base, let $O$ be open and let
  $x\in O$. Choose $\varepsilon>0$ such that $B(x,\varepsilon)\subseteq O$,
  and then choose $n$ with $2/n<\varepsilon$. Since the selected balls of
  radius $1/n$ cover $X$, there is $k$ such that
  \[
    x\in B(x_{n,k},1/n).
  \]
  For $y\in B(x_{n,k},1/n)$, the triangle inequality gives
  \[
    d(x,y)<\frac1n+\frac1n=\frac2n<\varepsilon.
  \]
  Hence
  \[
    x\in B(x_{n,k},1/n)\subseteq B(x,\varepsilon)\subseteq O,
  \]
  proving that $\mathcal{B}$ is a countable base.

  \textbf{(i) $\Rightarrow$ (iii).}
  Let $\mathcal{B}$ be a countable base. For every nonempty
  $B\in\mathcal{B}$, choose a point $x_B\in B$, and put
  \[
    D=\{x_B:B\in\mathcal{B},\ B\ne\varnothing\}.
  \]
  The set $D$ is countable. If $O$ is a nonempty open set, then it contains
  a nonempty basic open set $B\in\mathcal{B}$, so $x_B\in D\cap O$.
  Thus $D$ is dense, and $X$ is separable.

  \textbf{(iii) $\Rightarrow$ (i).}
  Let $D\subseteq X$ be a countable dense set. The family
  \[
    \mathcal{C}=\{B(a,1/n):a\in D,\ n\in\N\}
  \]
  is countable. Let $O$ be open and $x\in O$. Choose $\varepsilon>0$ such
  that $B(x,\varepsilon)\subseteq O$, and choose $n$ with $2/n<\varepsilon$.
  Since $D$ is dense, there is
  \[
    a\in D\cap B(x,1/n).
  \]
  Then $x\in B(a,1/n)$; moreover, if $y\in B(a,1/n)$, then
  \[
    d(x,y)\le d(x,a)+d(a,y)<\frac2n<\varepsilon.
  \]
  Consequently,
  \[
    x\in B(a,1/n)\subseteq B(x,\varepsilon)\subseteq O.
  \]
  So $\mathcal{C}$ is a countable base for $X$.
\end{solution}

\section{Question 3}

Let $(X,d)$ be a pseudometric space. Let $\mathcal{T}_d$ be the topology on
$X$ generated by
\[
  \{B(x,\varepsilon):x\in X,\ \varepsilon\in\R_{>0}\},
\]
where $B(x,\varepsilon):=\{y\in X:d(x,y)<\varepsilon\}$. Show that
$(X,\mathcal{T}_d)$ is $T_0$ if and only if $d$ is a metric.

\begin{solution}
  Suppose first that $d$ is a metric. If $x\ne y$, then $d(x,y)>0$.
  The open ball
  \[
    B\left(x,\frac{d(x,y)}2\right)
  \]
  contains $x$ but not $y$. Hence $(X,\mathcal{T}_d)$ is $T_0$.

  Conversely, suppose $(X,\mathcal{T}_d)$ is $T_0$. Let $x,y\in X$ satisfy
  $d(x,y)=0$. For every $\varepsilon>0$,
  \[
    y\in B(x,\varepsilon)\qquad\text{and}\qquad
    x\in B(y,\varepsilon).
  \]
  Therefore, every open set containing $x$ also contains $y$, and every
  open set containing $y$ also contains $x$. The $T_0$ property now implies
  $x=y$. Thus
  \[
    d(x,y)=0\implies x=y.
  \]
  A pseudometric already satisfies all the other metric axioms, so $d$ is a
  metric.
\end{solution}

\section{Question 4}

Prove the following for a topological space $(X,\mathcal{T})$:
\begin{enumerate}[(i)]
  \item if $X$ is $T_3$, then for every distinct $x,y\in X$, there exist
    open sets $U,V$ such that $x\in U$, $y\in V$, and
    $\overline{U}\cap\overline{V}=\varnothing$;
  \item if $X$ is $T_4$, then for every disjoint closed sets $A,B\subseteq X$,
    there exist open sets $U,V$ such that $A\subseteq U$, $B\subseteq V$, and
    $\overline{U}\cap\overline{V}=\varnothing$.
\end{enumerate}

\begin{solution}
  We first record two useful consequences of regularity and normality.

  If $X$ is regular, $x\in O$, and $O$ is open, then there is an open set
  $W$ such that
  \[
    x\in W\subseteq\overline{W}\subseteq O.
  \]
  Indeed, separate $x$ and the closed set $X\setminus O$ by disjoint open
  sets $W$ and $G$, respectively. Since $W\cap G=\varnothing$ and $G$ is
  open, $\overline{W}\subseteq X\setminus G\subseteq O$.

  Similarly, if $X$ is normal, $C\subseteq O$, $C$ is closed, and $O$ is
  open, then there is an open set $W$ such that
  \[
    C\subseteq W\subseteq\overline{W}\subseteq O.
  \]
  Separate the disjoint closed sets $C$ and $X\setminus O$ by disjoint open
  sets $W$ and $G$, respectively. The same closure argument gives
  $\overline{W}\subseteq X\setminus G\subseteq O$.

  \textbf{(i)} Since a $T_3$ space is $T_1$, the singleton $\{y\}$ is
  closed. By regularity, there are disjoint open sets $O_x,O_y$ such that
  \[
    x\in O_x,\qquad y\in O_y.
  \]
  Apply the regularity consequence above to $x\in O_x$ and $y\in O_y$.
  We obtain open sets $U,V$ satisfying
  \[
    x\in U\subseteq\overline{U}\subseteq O_x,
    \qquad
    y\in V\subseteq\overline{V}\subseteq O_y.
  \]
  Since $O_x\cap O_y=\varnothing$, it follows that
  \[
    \overline{U}\cap\overline{V}=\varnothing.
  \]

  \textbf{(ii)} Since $X$ is $T_4$, it is normal. Thus the disjoint closed
  sets $A$ and $B$ have disjoint open neighborhoods $O_A$ and $O_B$:
  \[
    A\subseteq O_A,\qquad B\subseteq O_B,\qquad O_A\cap O_B=\varnothing.
  \]
  Apply the normality consequence to the closed set $A$ contained in $O_A$
  and to the closed set $B$ contained in $O_B$. This gives open sets $U,V$
  with
  \[
    A\subseteq U\subseteq\overline{U}\subseteq O_A,
    \qquad
    B\subseteq V\subseteq\overline{V}\subseteq O_B.
  \]
  Therefore $\overline{U}\cap\overline{V}=\varnothing$.
\end{solution}

\section{Question 5}

Show that every locally compact Hausdorff space is regular, i.e.,
``local compactness $+\ T_2\Rightarrow T_3$.''

\begin{solution}
  Let $X$ be locally compact and Hausdorff. It is $T_1$ because every
  Hausdorff space is $T_1$. It remains to separate a point from a disjoint
  closed set.

  Let $x\in X$ and let $F\subseteq X$ be closed with $x\notin F$.
  By local compactness, $x$ has an open neighborhood $W$ contained in a
  compact set $K$. Since $X$ is Hausdorff, $K$ is closed. Put
  \[
    A=K\cap F.
  \]
  The set $A$ is closed in the compact space $K$, hence compact.

  For every $a\in A$, the Hausdorff property gives disjoint open sets
  $U_a,V_a\subseteq X$ such that
  \[
    x\in U_a,\qquad a\in V_a.
  \]
  The family $\{V_a:a\in A\}$ covers $A$. If $A\ne\varnothing$,
  compactness gives $a_1,\ldots,a_m\in A$ such that
  \[
    A\subseteq V_{a_1}\cup\cdots\cup V_{a_m}.
  \]
  If $A=\varnothing$, put $m=0$. In either case, define
  \[
    U=W\cap\bigcap_{j=1}^m U_{a_j},
  \]
  where the empty intersection is $X$.
  Then $U$ is an open neighborhood of $x$ and $U\subseteq K$. Hence
  $\overline{U}\subseteq K$, because $K$ is closed. Moreover,
  $U\cap V_{a_j}=\varnothing$ for every $j$, so, as $V_{a_j}$ is open,
  \[
    \overline{U}\cap V_{a_j}=\varnothing\qquad (1\le j\le m).
  \]
  Thus $\overline{U}\cap A=\varnothing$. Since $\overline{U}\subseteq K$,
  we obtain
  \[
    \overline{U}\cap F
    =\overline{U}\cap K\cap F
    =\overline{U}\cap A
    =\varnothing.
  \]
  Therefore $V=X\setminus\overline{U}$ is an open set containing $F$, and
  $U\cap V=\varnothing$. We have separated the arbitrary point $x$ from
  the arbitrary disjoint closed set $F$. Consequently $X$ is regular and,
  together with $T_1$, is a $T_3$ space.
\end{solution}

\end{document}
